\section{\texorpdfstring{\textbf{Q1}}{Q1} - A biblioteca proposta apresenta erro de quantificação equivalente ao das bibliotecas de referência para um mesmo método?}
\label{subsec:res_q1}

São comparados 28 métodos e 132 conjuntos de dados, distribuídos em 39 pares de método
e biblioteca. Cada par é avaliado sobre os conjuntos concluídos nas duas
implementações, de 108 a 132 conforme o caso.

A equivalência exige um teste próprio. A ausência de diferença significativa indica
apenas falta de evidência, de modo que afirmar igualdade pede a hipótese invertida. Foi
adotado o TOST sobre os postos sinalizados de Wilcoxon, com margem de $0{,}01$ em AE,
equivalente a um ponto percentual de prevalência. Abaixo desse valor a diferença não
altera decisão alguma tomada a partir da estimativa. A unidade de análise é o conjunto de
dados, com os quatro tamanhos de amostra reunidos pela média, uma vez que medidas do
mesmo conjunto não são independentes. A diferença central vem do estimador de
Hodges-Lehmann, acompanhada do intervalo de confiança de 95\% derivado do próprio teste,
e a multiplicidade dos 39 pares é corrigida por Holm.

O resultado é apresentado na Figura~\ref{fig:res_q1}, com um método contra uma
biblioteca de referência em cada linha. A diferença mediana de AE entre a
\texttt{mlquantify} e a referência é indicada pelo ponto, o intervalo de confiança pela
barra e a margem pelas linhas tracejadas. O intervalo contido na faixa é identificado
pelo marcador cheio, situação em que a compatibilidade está demonstrada. A biblioteca é
distinguida pela forma do marcador.

\begin{figure}[htb]
    \centering
    \includegraphics[width=0.86\linewidth]{Imagens/cap5/q1-equivalencia.pdf}
    \caption{Teste de equivalência do erro. O ponto é a diferença mediana de AE e a barra é o intervalo de confiança de 95\%.}
    \label{fig:res_q1}
    {\footnotesize Fonte: elaborado pelo autor.}
\end{figure}

Em 38 dos 39 pares o intervalo cabe na margem. Em 28 deles não há diferença detectável,
e nos 10 restantes existe diferença sistemática de magnitude irrelevante, sempre abaixo
de um ponto percentual. A maior delas alcança $0{,}0051$ em AE.

O único par fora da margem é o MS2 contra a \verb'QuaPy', com diferença mediana de
$-0{,}0078$ e intervalo de $[-0{,}0109,\,-0{,}0045]$, favorável à biblioteca proposta. A
divergência vem do conjunto de limiares candidatos. A \verb'QuaPy' percorre todos os
escores distintos obtidos na validação cruzada, enquanto a implementação proposta
percorre uma grade fixa de 101 pontos, o que desloca a mediana das estimativas tomadas
sobre esses candidatos.

Em RAE a diferença mediana entre os pares fica em $0{,}0001$ e alcança $0{,}0329$ no
extremo, com intervalo contendo zero em 20 dos 39 pares. A margem de $0{,}01$ está na
unidade da prevalência e vale apenas para o AE, razão pela qual o RAE é reportado sem
veredito de equivalência.

A equivalência em erro fica sustentada por esse resultado. Nas implementações da
biblioteca proposta, o comportamento das referências é reproduzido dentro de uma margem
sem consequência prática, e a única exceção favorece a biblioteca.


\section{\texorpdfstring{\textbf{Q2}}{Q2} - Qual é o custo em tempo e memória da biblioteca proposta no treinamento e na estimação quando comparada com as demais bibliotecas?}
\label{subsec:res_q2}

O custo é comparado pela razão entre as medidas, tratada em logaritmo de base 2, única
escala em que o dobro e a metade ficam à mesma distância da igualdade. O ponto é a razão
mediana entre conjuntos pelo estimador de Hodges-Lehmann e a barra é o intervalo de 95\%,
na mesma gramática da figura anterior. O tempo soma as fases de treinamento e estimação,
que são trabalho sequencial. A memória toma o maior pico entre as duas, uma vez que o
processo libera a primeira fase antes de iniciar a segunda.

O tempo é apresentado na Figura~\ref{fig:res_q2_tempo}. A biblioteca proposta é mais barata com
intervalo inteiro abaixo de 1 em 18 dos 39 pares e mais cara em 14, com razão mediana de
$0{,}967$. A vantagem se concentra nos métodos de casamento de distribuição multiclasse,
com razão de $0{,}016$ no GHDx e de $0{,}023$ no GACC, ambos contra a \verb'qunfold'. A desvantagem aparece na contagem ajustada da
\verb'QuantificationLib', com $3{,}4$ no PWKAC e no PACC e $3{,}0$ no FM, onde
aquela biblioteca resolve o ajuste em forma fechada e a proposta recorre a otimização com
restrição a cada estimativa.

\begin{figure}[htb]
    \centering
    \includegraphics[width=0.78\linewidth]{Imagens/cap5/q2-tempo.pdf}
    \caption{Razão do tempo total da \texttt{mlquantify} em relação a cada biblioteca de referência, somadas as fases de treinamento e estimação.}
    \label{fig:res_q2_tempo}
    {\footnotesize Fonte: elaborado pelo autor.}
\end{figure}

A memória é apresentada na Figura~\ref{fig:res_q2_memoria}, com vantagem mais uniforme.
A biblioteca proposta ocupa menos em 32 dos 39 pares e mais em 6, com razão mediana de
$0{,}816$. Os extremos vão de $0{,}19$ no GHDx contra a \verb'qunfold' a $1{,}49$ no
KDEyCS contra a \verb'QuaPy'.

\begin{figure}[htb]
    \centering
    \includegraphics[width=0.78\linewidth]{Imagens/cap5/q2-memoria.pdf}
    \caption{Razão do pico de memória da \texttt{mlquantify} em relação a cada biblioteca de referência, tomado o maior valor entre as duas fases.}
    \label{fig:res_q2_memoria}
    {\footnotesize Fonte: elaborado pelo autor.}
\end{figure}

As fases são separadas na Tabela~\ref{tab:res_q2_fases}, em que se observam
comportamentos distintos. Contra
a \verb'QuaPy' e a \verb'QuantificationLib' o treinamento da biblioteca proposta é
ligeiramente mais barato e a estimação é ligeiramente mais cara, com razões próximas da
igualdade nas quatro medidas. Contra a \verb'qunfold' a diferença muda de ordem de
grandeza, com $0{,}017$ no tempo de estimação, efeito do otimizador numérico empregado
por aquela biblioteca em cada estimativa.

\begin{table}[htb]
    \centering
    \caption{Razão mediana de custo da \texttt{mlquantify} em relação a cada biblioteca de referência, por fase. Valores abaixo de 1 indicam custo menor.}
    \label{tab:res_q2_fases}
    \footnotesize
    \renewcommand{\arraystretch}{1.25}
    \begin{tabular}{lccc}
        \hline
        \textbf{Medida} & \verb'QuantificationLib' & \verb'QuaPy' & \verb'qunfold' \\
        \hline
        Tempo de treinamento   & 0,907 & 0,972 & 1,080 \\
        Tempo de estimação     & 1,021 & 1,053 & 0,017 \\
        Memória no treinamento & 0,924 & 0,789 & 0,998 \\
        Memória na estimação   & 0,743 & 0,904 & 0,388 \\
        \hline
    \end{tabular}

    \vspace{2pt}
    {\footnotesize Fonte: elaborado pelo autor.}
\end{table}

O custo da biblioteca proposta fica na mesma ordem de grandeza das referências, com
vantagem clara em memória e equilíbrio em tempo. As perdas se concentram em métodos cuja
formulação admite solução fechada, e os ganhos nos métodos multiclasse de casamento de
distribuição.


\section{Considerações Finais}
\label{sec:consideracoes_finais_res}

Neste capítulo, foram apresentadas a \texttt{mlquantify} e a comparação com as
bibliotecas de referência. Na arquitetura, os métodos são reunidos sob uma interface
única, com \textit{mixins} e \textit{tags} para exigências distintas, decomposição
multiclasse automática e composição explícita de representação, perda e solucionador.

As duas questões formuladas foram respondidas pela avaliação experimental. Em erro, o
intervalo de confiança ficou dentro da margem de um ponto percentual de prevalência em
38 dos 39 pares de método e biblioteca, e o único par fora dela favorece a biblioteca
proposta. Em custo, a razão mediana ficou em $0{,}967$ no tempo e em $0{,}816$ na
memória, com ganhos nos métodos multiclasse de casamento de distribuição e perdas nos
métodos de contagem ajustada que a \verb'QuantificationLib' resolve em forma fechada.

No próximo capítulo, são discutidas as limitações do trabalho e as possibilidades de
continuidade.
